\subsection{実務での使い分け}
モデルや方法の選択で最も大切なのは、目的に合っているかである。すべてのプロジェクトで同じモデルを作る必要はない。逆に、必要なモデルを作らないと、後で誤解、欠陥、変更コストとして戻ってくる。

\par\smallskip\noindent\refstepcounter{table}\label{tab:ch11-09}表 \thetable: 実務での使い分けの整理\par\smallskip
表 \thetable は、実務での使い分けの整理を比較・確認しやすい形で整理したものである。\par\smallskip
\begin{longtable}{>{\raggedright\arraybackslash}p{0.31\textwidth}>{\raggedright\arraybackslash}p{0.64\textwidth}}
\toprule
状況 & 向きやすい選択 \\
\midrule
業務データが複雑で、関係の誤りが重大である & 情報モデル、概念データモデル、構造モデルを重視する。\\
状態遷移や例外処理が多い & 状態機械、活動図、シーケンス図などの振る舞いモデルを重視する。\\
安全性やセキュリティが重要で、曖昧さを減らしたい & 形式手法、軽量形式手法、モデル検査、事前条件・事後条件の記述を検討する。\\
利用者が要求を言語化しにくい & プロトタイピング、ユーザインタフェース試作、シナリオ確認を使う。\\
要求が変わりやすく、顧客から頻繁に学ぶ必要がある & アジャイル方法、短い反復、継続的フィードバックを使う。\\
規制、契約、安全認証が強い & 計画駆動方法、追跡可能性、形式的なレビュー、記録の管理を重視する。\\
\bottomrule
\end{longtable}

\begin{pointbox}{実務上の判断基準}
モデルや方法は、作る負担だけで評価してはいけない。誤解を減らす効果、欠陥を早期に見つける効果、関係者の合意を助ける効果、変更時の影響分析を容易にする効果も含めて評価する必要がある。
\end{pointbox}
